%--------------------------------------------------------------
\section{Graph Transformers}
\label{sec:gps}
%--------------------------------------------------------------

Graph Transformers leverage the transformers self-attention mechanism to enrich the local message passing mechanism with global context awareness. 
In this work we adopt the \emph{General, Powerful, Scalable} (GPS) architecture ~\cite{rampasek2022gps}, implemented in PyTorch Geometric as \texttt{GPSConv}.
The main idea is to combine, within each layer, local message-passing toghether with global self-attention, allowing the model to capture both local structural information and long-range dependencies.

\subsection{GPSConv Layer}
In each GPSConv layer the node embeddings are updated by two parallel procedures, local message passing and global attention; the two outputs are then combined to produce the final output. 

\paragraph{Local Message Passing}
The local message passing is performed by a Graph Isomorphic Network (GIN), as the one presented in Section~\ref{sec:gin}.
\begin{equation}
  h_i^{\text{local},(k)} = \operatorname{MLP}_{\text{local}}^{(k)}\!\left(
    \left(1 + \varepsilon^{(k)}\right) h_i^{(k-1)}
    + \sum_{j \in \mathcal{N}(i)} h_j^{(k-1)}
  \right)
\end{equation}

\paragraph{Global Attention}
A standard multi-head self-attention module is applied over all nodes, capturing long-range interactions without the locality constraint of message passing:
\begin{equation}
  h_i^{\text{global},(k)} = \operatorname{MHA}^{(k)}\!\left(
    h_i^{(k-1)},\;\left\{h_j^{(k-1)}\right\}_{j \in V}
  \right)
\end{equation}
with $H$ parallel attention heads.
\bigskip

The outputs of both branches are then summed toghether to obtain the final node representation:
\begin{equation}
  h_i^{(k)} = h_i^{\text{local},(k)} + h_i^{\text{global},(k)}
\end{equation}

\subsection{Input/Output Projections}

GPSConv operates with a fixed channel dimension $d_h$ across all layers.
Since the original node features $x_i \in \mathbb{R}^{d_{\text{feat}}}$ and
the desired output dimension $d_{\text{out}}$ generally differ from $d_h$,
two linear projections are added:
\begin{align}
  \tilde{x}_i &= W_{\text{in}}\, x_i, \quad W_{\text{in}} \in \mathbb{R}^{d_h \times d_{\text{feat}}} \\
  X^{(L)} &= W_{\text{out}}\, H^{(L)}, \quad W_{\text{out}} \in \mathbb{R}^{d_{\text{out}} \times d_h}
\end{align}
where $H^{(L)} \in \mathbb{R}^{N \times d_h}$ is the matrix containing the final node embeddings $h_i^{(L)}$ for nodes after the $L$-th layer.
Notice that the constraint $d_h \bmod H = 0$ must be satisfied to allow MHA to split the embedding evenly across attention heads.

The intermediate representations produced by each head are then concatenated along the feature dimension and linearly projected to form the final multi-head attention output.

By combining local GIN message passing with global multi-head self-attention, GPSConv overcomes the WL ceiling and oversquashing.
